\chapter*{Abstract}
\addcontentsline{toc}{chapter}{Abstract}

A prediction market must expose a participant's belief twice: once to whoever prices the
position, and once to whoever decides what happened. Where the traders are agents and the
resolver is a language model, both exposures change character, and two protocols now running
make opposite choices at each. Delphi matches against a dynamic pari-mutuel pool, adding no
pre-trade disclosure of its own, and settles by AI judgement of a specification fixed in
advance. Sapience broadcasts a request for quotes before any fill and resolves through a
bonded optimistic-oracle assertion. This report asks where strategic informational advantage
sits in such markets, what it is worth, and whether it is AI-MEV: value captured through
role-specific protocol access, beyond what forecasting ability or ordinary compensation for
risk would earn.

From published documentation and implementation source, the report establishes that a
Sapience request reaches every connected client, unfiltered, carrying the direction of every
leg, the position size and the requester's address, with no way to disclose less and still
trade. A simulator measures what that costs against a matched benchmark in which responders
cannot condition on the request. Disclosure removes about seventy-three per cent of an
informed forecaster's expected surplus at the baseline; direction alone accounts for roughly
ninety-five per cent of that, so coarsening size is nearly useless as a mitigation; and the
share lost rises with the forecaster's accuracy, to about ninety-six per cent for the
best-informed forecaster modelled.

The report does not settle whether that transfer is AI-MEV. Of the definition's two limbs,
excess over ordinary compensation is satisfied on the protocol's own documentation, while
role-specificity turns on whether capital-gated access to a public stream counts as
protocol-specific, a question the MEV literature has not resolved. What survives regardless
is a design critique: the mechanism compels complete, address-linked disclosure, offers
nothing cheaper, and the entire measured transfer accrues to the responders who quote.

The second contribution needs no experiment. Semantic validity, reproducibility and
execution attestation are three distinct properties; the latter two are not substitutes,
because commercial models are non-deterministic; and no verification mechanism guarantees
the first. Both protocols have built serious machinery against dishonesty. What survives it in
each case is a disclosure the mechanism compels and does not price, and an ambiguity it
cannot adjudicate; neither is dishonesty, which is why neither is caught. The model cannot separate information rent from risk
compensation, its benchmark is not a market any responder would join, and no live venue was
measured.
